\begin{abstract}

Domain-specific systems-on-chip (SoCs) can substantially improve application performance and energy efficiency, but their design still requires extensive cross-layer expertise and costly manual effort. Existing SoC frameworks automate integration and verification, yet application analysis, hardware/software partitioning, accelerator generation, and system-level optimization remain largely designer-driven.
We present SoCPilot, an agentic design framework that leverages large language models (LLMs) for end-to-end SoC customization. Starting from a C/C++ application and design constraints, SoCPilot combines dynamic profiling and static analysis to identify computational hotspots, then performs hardware/software partitioning to select functions for acceleration. For the selected hotspot functions, LLMs insert HLS pragmas and generate synthesizable accelerators. Afterwards, SoCPilot integrates these accelerators with a selected processor and memory-system configuration and performs system-level design space exploration over latency and estimated area. The generated SoC is verified through simulation and validated on an FPGA board.
Evaluations on four FPGA-validated applications show that the generated SoCs achieve a geometric mean end-to-end speedup of 2.66$\times$ over matched CPU-only execution. These results demonstrate that LLMs can coordinate application analysis, hardware generation, tool execution, and cross-layer optimization to support application-driven SoC customization.

% Modern system-on-chip (SoC) design faces escalating challenges, including increasing functional complexity, verification bottlenecks, architectural demands, and shrinking time-to-market windows, all of which require substantial human effort and result in low productivity. Although many design frameworks have been developed to mitigate these issues through automated integration and verification, they often support limited configurations and still demand extensive expertise in both hardware and software.
% Leveraging the recent advances of large language models (LLMs) in code generation, tool usage, and reasoning, we propose SoCPilot, an LLM-driven framework that automates the end-to-end System-on-Chip (SoC) design flow. SoCPilot empowers LLMs to interpret user requirements, select appropriate IPs via semantic search, and construct task dependency graphs. LLMs also collaborate with simulators and PPA predictors to estimate performance and guide HW/SW partitioning through heuristic optimization. 
% For hardware-mapped tasks, LLMs generate specialized accelerators from high-level descriptions and orchestrate their integration with other IPs using the ESP platform that automates SoC construction and physical design. The framework further conducts iterative design space exploration and validates the design on FPGAs. By embedding LLMs into the SoC design flow, SoCPilot significantly lowers design complexity and enhances productivity of the SoC development.
% Experimental results show that SoCPilot achieves fully automated SoC design from natural language and high-level applications to SoC layouts. The end-to-end process completes within 24 hours and delivers an average performance speedup of 6.9x compared to pure software solutions, demonstrating the great potential of LLM-assisted SoC design automation.


 
\end{abstract}

\begin{IEEEkeywords}
	large language model, high-level synthesis, system-on-chip, hardware/software partitioning
\end{IEEEkeywords}


